\section{Kiểm thử và đánh giá}

Kế hoạch kiểm thử của dự án tập trung vào ba mức: unit test để kiểm tra từng đơn vị nghiệp vụ quan trọng, E2E test để xác nhận các luồng API chạy qua NestJS application và PostgreSQL, và CI để tự động kiểm tra chất lượng khi có thay đổi mã nguồn. Các nhóm kiểm thử ưu tiên những rủi ro chính của hệ thống Feature Flag: evaluation phải xác định, mặc định an toàn, mutation phải có audit log, cache không làm sai kết quả, phân quyền phải được xử lý ở backend và SDK phải fail-closed khi lỗi.

\subsection{Unit test}

Unit test được triển khai bằng Jest trong hai workspace chính: \texttt{@ffp/backend} và \texttt{@ffp/js-sdk}. Ở backend, unit test kiểm tra Rule Evaluation Engine, stable rollout hash, evaluation service/controller/repository, cache, invalidation, audit log, history, RBAC, DTO validation, exception filter, request context, repository và statistics service. Các test dùng mock hoặc dữ liệu tối thiểu để cô lập đơn vị cần kiểm tra, sau đó đối chiếu output với contract mong đợi như \texttt{enabled}, \texttt{reason}, error code, audit snapshot hoặc permission result.

Với JavaScript SDK, unit test kiểm tra client gọi Evaluation API, validate response, trả helper \texttt{isEnabled}/\texttt{getVariant}, xử lý timeout, HTTP error, response sai định dạng và fallback fail-closed. Nhóm test này bảo đảm ứng dụng khách không tự bật tính năng khi Evaluation API không khả dụng.

\begin{table}[H]
    \centering
    \small
    \caption{Thống kê unit test đã triển khai}
    \label{tab:unit-test-summary}
    \begin{tabular}{p{0.24\textwidth} p{0.36\textwidth} p{0.16\textwidth} p{0.14\textwidth}}
        \toprule
        \textbf{Workspace} & \textbf{Phạm vi kiểm thử chính} & \textbf{Test suite} & \textbf{Test case} \\
        \midrule
        \texttt{@ffp/backend} & Evaluation engine, rollout hash, cache, audit, RBAC, DTO, controller, service, repository, statistics & 52 & 416 \\
        \texttt{@ffp/js-sdk} & SDK client, response validation, helper API, timeout, lỗi mạng và fallback fail-closed & 1 & 21 \\
        Tổng & Unit test toàn workspace qua \texttt{npm run test -- --runInBand} & 53 & 437 \\
        \bottomrule
    \end{tabular}
\end{table}

\subsection{Integration test và E2E test}

Repository đã cấu hình Jest cho integration test qua \texttt{apps/backend/test/jest-integration.json} và script \texttt{npm run test:integration --workspace=@ffp/backend}. Ở trạng thái hiện tại, không có file \texttt{*.integration-spec.ts} riêng; các kiểm thử tích hợp được thể hiện chủ yếu qua nhóm E2E test trong \texttt{apps/backend/test}.

E2E test dùng \texttt{@nestjs/testing} để khởi tạo \texttt{AppModule}, áp dụng global validation pipe, exception filter, request context middleware và demo RBAC giống runtime. Các test gọi HTTP bằng \texttt{supertest}, thao tác dữ liệu qua Prisma và dọn dữ liệu bằng helper test database. Hiện có 10 file \texttt{*.e2e-spec.ts}, bao phủ health check, management API, vertical slice tạo project/flag/rule/evaluate, API hardening, demo flow, flag history, group kill switch, evaluation cache, evaluation statistics và server-resolved RBAC.

\subsection{Luồng CI trong CI/CD}

Luồng CI được triển khai bằng GitHub Actions tại \texttt{.github/workflows/ci.yml}. Workflow chạy khi push hoặc pull request vào nhánh \texttt{develop} và \texttt{main}, đồng thời hỗ trợ chạy thủ công bằng \texttt{workflow\_dispatch}. Job \texttt{validate} sử dụng Ubuntu 24.04, Node.js 22.22.2 và PostgreSQL 16 service để mô phỏng môi trường kiểm thử có cơ sở dữ liệu thật.

Các bước CI lần lượt gồm kiểm tra whitespace bằng \texttt{git diff --check}, cài dependencies bằng \texttt{npm ci}, validate và generate Prisma client, chạy lint, chạy unit test backend, build toàn bộ workspace, apply database migration, chạy seed hai lần để kiểm tra idempotency, chạy integration test và cuối cùng chạy backend E2E test. Luồng này giúp phát hiện sớm lỗi định dạng, lỗi schema, lỗi TypeScript/build, lỗi unit test và lỗi end-to-end trước khi thay đổi được hợp nhất.
